\begin{lemma}[Common exact-edge and cutoff preparation for the disjoint branch]
Fix the graph, palette, distinguished edge, input, exact terminal tree,
generating coloring, positive-rung ETT, hierarchy base, counterexample
\(C\), minimum context, and replacement presentation \(P\).
Write \(\mu=P.\mathit{coloring}\), \(H=P.\mathit{counterexample.hierarchy}\),
\(g=P.\mathit{counterexample.good}\), \(q=H.\mathit{dividerCount}\),
\(j=\noderef{HierarchyPathNumber}(C.\mathit{hierarchy})\),
\(\ell=\noderef{HierarchyFullTailLength}(C)\),
\(A=\noderef{HierarchyStarredStage}(H,q)\),
\(B=\noderef{HierarchyFullTailPrefix}(C,j-2)\),
\(X=\noderef{HierarchyFullTailPrefix}(C,j-1)\).
Let \(T_n\) be the terminal ladder tree of the replacement ETT,
\(D=\noderef{HierarchyStoredColorsAt}(H,q)\),
\(\Gamma=\bigcup_{\eta\in D}g.\Gamma^q_\eta\), and
\[
 M_0=\bigcup_{v\in V(P.\mathit{replacementBase.core})\setminus V(T_n)}
       \noderef{MissingColors}(G,\mu,v),\qquad
 U_j=\{c:\exists r\in\operatorname{Fin}|T|,+ |A|\le r<|A|+j,\ \mu(T.\mathit{edge}(r))=c\}.
\]
Assume \(2\le j<\ell\), \(\alpha\) is missing at \(y_\ell\) and on
\(B\), and either \(\alpha\notin\Gamma\cup M_0\) or there is
\(\eta\in D\) with \(\alpha\in g.\Gamma^q_\eta\) and first
\(\eta\)-missing position on \(T\) at most \(|B|\).
Assume also
\(\beta\ne\alpha\), \(\beta\in\overline\mu(B)\),
\(\beta\notin\Gamma\cup U_j\), and either
\(\beta\notin M_0\cup D\), or
\(\beta\in D\) and \((g.\Gamma^q_\beta\cup\{\beta\})\cap U_j=\varnothing\).
The component \(Q=\noderef{TwoColorComponent}(G,\mu,\alpha,\beta,y_\ell)\)
contains no vertex-node of \(X\).
Then, for \(\sigma=\noderef{TwoColorComponentSwap}(G,\mu,\alpha,\beta,y_\ell)\),
there exists a \noderef{P2880_T1_SplitSwapData}.
\end{lemma}
\begin{proof}
All properness, exact-one-uncoloredness, maximum-property and induction
data below belong to \(P\), by \noderef{HierarchyReplacementPresentation}.
\noderef{P2752_T2_HierarchyMeasures} applied to its same-good stage
relation identifies the old and replacement starred stage. Expanding
\noderef{HierarchyFullTailLength} and \noderef{HierarchyFullTailPrefix}
therefore identifies all their numerical full-tail cutoffs and exact
prefixes. In particular the replacement penultimate prefix is the exact
tree \(Z\) of length \(|T|-1\).
\noderef{HierarchyCounterexamplePenultimateElementary}, applied to all
six fields of \(P.\mathit{minimumContext}\), makes \(Z\) elementary
under \(\mu\).

The hierarchy's nested stages and core-inside-first field in
\noderef{TashkinovHierarchy} place the base core in \(A\) and
\(A\) in \(T\). At zero dividers \(A\) is the core; at positive
dividers it is the last retained ordinary stage. Exact prefixes of a
common sequence compare by their lengths: their vertex and edge entries
are the corresponding entries of that common sequence. Since
\(2\le j<|T|-|A|\), the minimum cutoffs in the full-tail prefixes
are unsaturated, and their lengths are
\[
 |B|=|A|+j-2,\quad |X|=|A|+j-1,\quad |Z|=|T|-1.
\]
Consequently \(A\preccurlyeq B\prec X\preccurlyeq Z\prec T\);
the core is a prefix of \(X\) and \(|A|<|X|\).
The equality \(B=A\) at \(j=2\) is permitted. All these comparisons
are exact \noderef{TreeSequenceInitialSegment} relations.

The root is a vertex of its own connected component. Because \(\alpha\)
is missing there, all incident two-color edges have color \(\beta\).
Properness allows at most one such edge identity, so its incidence
neighbor set is subsingleton. Together with \(\alpha\ne\beta\),
this is precisely \noderef{TwoColorPathComponent}, also when the root
is isolated. \noderef{TwoColorComponentSwapProper} proves that
\(\sigma\) is proper. The formula of \noderef{TwoColorComponentSwap}
fixes \(\mathsf{none}\), interchanges only the two colored values,
and fixes all other values. Its uncolored edge set is therefore the same
singleton \(\{e_0\}\).

Apply \noderef{P2880_T1_ElementaryPrefixSwap} to \(B\preccurlyeq
X\preccurlyeq Z\), the two missing witnesses, and the assumed
component exclusion. It fixes every edge incident with \(X\) and
identifies all missing sets on \(Z\). Each displayed edge below
\(|X|\) has an endpoint in \(X\), so it is fixed.
This gives \noderef{TreeColoringInvariant} on \(X\).
The one further edge with index \(|X|=|A|+j-1\) has its attaching
endpoint at a preceding vertex of \(X\), by the attachment field of
\noderef{TreeSequence}. It too is incident with \(X\) and is fixed.
Thus every individual edge counted by \(U_j\) is fixed.
\noderef{TreeSequenceInitialSegmentVertexSetSubset} and equality of the
missing sets transport the two old witnesses in \(B\) to \(\sigma\)
on \(B\) and on \(X\), and retain them under \(\mu\) on \(Z\).
The root-membership and affected-missing clause of
\noderef{TwoColorComponentSwapMissingBehavior} show that \(\beta\)
is missing at \(y_\ell\) under \(\sigma\), without assuming that it
was present there beforehand.

For \(c\in\{\alpha,\beta\}\), choose its old missing witness in
\(B\). The finite minimum in
\noderef{HierarchyFirstMissingPosition} puts its first missing vertex
on \(T\) at index at most \(|B|\), hence in \(X\). That vertex-node
is excluded from \(Q\), so its component differs from \(Q\).
\noderef{TwoColorComponentsDisjoint} gives the two claimed component
separations, separately, even if these first missing vertices coincide.

Fix any \(\eta\in D\). If \(\eta\) is neither affected color,
\noderef{TwoColorComponentSwapMissingBehavior} identifies its missing
predicate at every vertex of \(T\). The finite candidate sets for its
first position, and their terminal fallback when empty, are identical.
If \(\eta\) is one of the affected colors, both old and new predicates
have a witness in \(B\). Their first positions are at most \(|B|\).
Their vertexwise missing predicates on \(B\) agree, and
\noderef{HierarchyFirstMissingPositionPrefixCutoff} identifies the
truncated first position on \(B\) with that on \(T\) for each
coloring. Since each first position is already at most \(|B|\), the
truncations yield equality of the full first positions.
For the affected-Gamma cutoff, \(\beta\notin\Gamma\) excludes that
membership. If \(\alpha\notin\Gamma\), the implication is vacuous.
Otherwise its supplied owner \(\eta_0\in D\) has first position
at most \(|B|\). Pairwise Gamma disjointness in
\noderef{GoodHierarchy} implies that any other owner of \(\alpha\)
is \(\eta_0\). Hence every affected Gamma owner has first position
at most \(|B|<|X|\).

For the augmentation routes use \noderef{P2813_T1_TerminalRoute} on
the actual replacement ETT, not a hypothetical terminal constructor.
It gives either the PE route from \(T_n\) to \(T\), or the SE/RE
route from the actual seed \(S=T_n+f_n\), with the exact singleton
connecting-edge set and \noderef{TreeSequenceAdjoinEdge} certificate.
The PE source is a prefix of the core, by the
\noderef{OrderedTreeUnion} in \noderef{HierarchyBase}, and so a
prefix of \(X\). In SE/RE the nonparallel core is \(T_n\).
Thus \(|S|=|T_n|+1\le |A|+1\le |X|\).
Both \(S\) and \(X\) are prefixes of \(T\), so \(S\preccurlyeq X\).
This includes \(S=X\) at zero dividers and \(j=2\); it never
asserts the false comparison \(S\preccurlyeq A\) there.
\noderef{TashkinovAugmentationSequenceRestrict} restricts the old
route to \(X\), and \noderef{TashkinovAugmentationSequenceSuffix}
gives the complementary old suffix \(X\to T\).
If the source equals \(X\), its new prefix route is reflexive.
Otherwise transport the restricted route by
\noderef{TashkinovAugmentationSequenceColoringTransport}, using the
exact penultimate prefix \(B\prec X\) of length \(|X|-1\),
all fixed edges below \(X\), and equality of the missing sets on
\(B\). This retains the exact SE/RE seed when present.
Finally \noderef{TashkinovAugmentationSequenceComponentSwap}, with
the two new missing witnesses on \(X\), transports the old suffix to
\(\sigma\). These arguments supply every field of
\noderef{P2880_T1_SplitSwapData}, whose swap-equality field is reflexive.
\end{proof}
